\section{Entropy enrichment from a forward segment}\label{sec:geometry}

Throughout this section, \(z_0,z_1,\ldots\) is a fixed self-avoiding walk
in \(\Z[i]\) with \(|z_{t+1}-z_t|\le D\), where \(D\ge1\).
The walk need not satisfy any congruence restrictions.
Our goal is a forward sampling operation that raises the average
entropy of residue projections. Two geometric estimates supply the
input: a walk segment has many distinct differences, and random signed
products usually separate those differences. We then sample an endpoint
of a uniformly chosen distinct difference. Its residue entropy inherits
both the entropy already present at the starting point and the entropy
created within the segment.

\subsection{How geometry produces many differences}

\begin{lemma}[Many differences]\label{lem:differences}
Let \(E\) consist of the points of a segment of \(n\ge1\) steps of the
walk. Let \(R\) be its diameter, choose a line through two points
realizing that diameter, and let \(W\) be the larger of \(1\) and the
maximum distance of a point of \(E\) from that line. Then, for \(A=RW\),
\[
 c n\le A\le D^2n^2,
 \qquad |E-E|\ge c_D A,
\]
where \(c>0\) is absolute and \(c_D>0\) depends only on \(D\).
\end{lemma}

\begin{proof}
Let \(p_0,p_1\) be the diameter endpoints. Every point of \(E\) projects
onto the segment \(p_0p_1\): a projection beyond either endpoint would
give a distance greater than \(R\) to the other endpoint.
Thus \(E\) lies in a rectangle of length \(R\) and width \(2W\).
Here \(R\ge W\ge1\). Disjoint disks of radius \(1/3\) about the distinct
lattice points of \(E\) lie in a fixed enlargement of this rectangle.
Its area is \(O(RW)\), proving \(n+1=O(A)\).
The step bound gives \(R\le Dn\) and \(W\le R\), proving the upper bound.

If \(W=1\), subtracting a fixed member of \(E\) gives
\(|E-E|\ge n+1\ge A/D\).
Suppose instead that \(W>1\), and choose \(p_2\in E\) at distance \(W\)
from the diameter line. Put
\[
 a=p_1-p_0,\qquad b=p_2-p_0,\qquad
 \Lambda=\Z a+\Z b.
\]
The flat torus \(\R^2/\Lambda\) has area
\(|\det(a,b)|=RW=A\).
Follow the interpolated walk, in either time direction as necessary,
to obtain paths \(\alpha,\gamma:[0,1]\to\R^2\) from \(p_0\) to \(p_1,p_2\),
respectively. The map
\[
 (s,t)\longmapsto \alpha(s)-\gamma(t)\pmod{\Lambda}
\]
descends to a continuous map from the parameter torus
\([0,1]^2/\!\sim\) to \(\R^2/\Lambda\): opposite edges agree because
the respective endpoint differences are \(a\) and \(-b\).
Homotoping both paths to their straight segments while keeping their
endpoints fixed shows that this map has degree \(-1\), with the target
oriented by \(a,b\). In particular it is onto.
Indeed, a map omitting a point factors through a punctured torus,
whose second homology is zero, and therefore has degree zero.
These are the standard homological properties of degree;
see~\cite[Section 3.3, Theorem 3.26(a), Proposition 3.29, and Exercise 7]{Hatcher}.

Every interpolated point is within distance \(D\) of a vertex of its
walk edge. Consequently the images of \(E-E\), with disks of radius
\(2D\) about them, cover the target torus.
Each such projected disk has area at most \(4\pi D^2\), regardless of
overlap or wrapping around the torus. Hence
\[
 A\le 4\pi D^2 |E-E|,
\]
as required.
\end{proof}

\subsection{Separating differences with randomly chosen factors}

Many distinct differences give high projected entropy only if reduction
separates them. The next lemma excludes collisions by ruling out nonzero
multiples in a rectangle containing differences of differences.
Recall that a \emph{sign choice} over a split rational prime means
selecting one of its two conjugate Gaussian factors, with equal probability.

\begin{lemma}[A randomly signed product avoids a thin rectangle]
\label{lem:rectangle}
Let \(p_1,\ldots,p_r\) be distinct rational primes congruent to \(1\)
modulo \(4\), all in \([T,2T]\), where \(T\ge5\), and put \(P=\prod_i p_i\).
Let \(Q\) be an origin-centered rectangle of half-lengths
\(R_1\ge W_1\ge1\), in any orientation.
Suppose
\[
 R_1W_1\le P^{1-d},\qquad 0<d<\tfrac14.
\]
There are absolute positive constants \(c,C,C_0\) such that, if
\(d^3r\ge C_0\), independent sign choices satisfy
\[
 \Pr\left\{Q\cap\Bigl(\prod_{i=1}^r\pi_i\Bigr)\Z[i]
                      \ne\{0\}\right\}
 \le C e^{-c d^3r}.
\]
\end{lemma}

\begin{proof}
Call a sign vector bad if it has a nonzero witness in \(Q\).
We group bad sign vectors by the lines containing their witnesses.
Nearby sign vectors will have the same line, whereas a fixed line will
occur for only a small fraction of signs. A cube expansion estimate
then combines these facts without counting the possible lines.

\paragraph{Nearby signs have collinear witnesses.}
We first show that witnesses belonging to nearby bad sign vectors lie
on a common line through zero.
If two sign vectors differ in \(h\le dr/4\) coordinates, the product
\(c_*\) of their common Gaussian factors has
\[
 N(c_*)\ge P/(2T)^h.
\]
For corresponding witnesses \(w=c_*u\) and \(w'=c_*v\),
\[
 \det(w,w')=N(c_*)\det(u,v)\in N(c_*)\Z.
\]
On the other hand, \(|\det(w,w')|\le2R_1W_1\le2P^{1-d}\).
The logarithm of the ratio between the divisor and this upper bound
is at least
\[
 dr\left(\log T-\tfrac14\log(2T)\right)-\log2,
\]
which is positive when \(d^3r\) exceeds a sufficiently large absolute
constant. The determinant must therefore vanish.
This also applies to two witnesses for the same sign vector, so every
bad sign vector has a unique witness line.

\paragraph{A fixed line is unlikely.}
Fix a line through zero with primitive integer direction
\(v\in\Z[i]\setminus\{0\}\), meaning that the two integer
coordinates of \(v\) are relatively prime.
Its lattice points are precisely the integer multiples \(a v\),
\(a\in\Z\). Define
\[
 P_\sigma(v)=\prod_{\pi_i\mid v}p_i
\]
for sign vector \(\sigma\). If \(\pi_i\nmid v\) and
\(\pi_i\mid a v\), then \(\pi_i\mid a\), and for the rational integer
\(a\) this is equivalent to \(p_i\mid a\).
Thus any nonzero witness on this line has length at least
\[
 |v|\,P/P_\sigma(v).
\]
Primitivity of the two integer coordinates of \(v\) implies that at most
one factor over any \(p_i\) divides \(v\). Moreover, the norm of the
product of all such eligible factors is at most \(N(v)=|v|^2\).
It follows that
\[
 \E_\sigma\log P_\sigma(v)\le\log|v|.
\]
Every point of \(Q\) has length at most \(\sqrt2R_1\), and
\(R_1\le P^{1-d}\). Existence of a witness on the fixed line therefore
requires
\[
 \log P_\sigma(v)\ge\log|v|+d\log P-\log\sqrt2.
\]
The random variable on the left is a sum of independent variables,
each of range at most \(\log(2T)\). Hoeffding's inequality
\cite[Theorem 2, Eq.~(2.6)]{Hoeffding} gives, after increasing \(C_0\) if needed,
\begin{equation}\label{eq:line-tail}
 \Pr_\sigma\{\text{a witness lies on the fixed line}\}
 \le \exp(-c_0d^2r)
\end{equation}
for an absolute \(c_0>0\). Here we used
\(\log P\ge r\log T\) and the bounded ratio
\(\log(2T)/\log T\).

\paragraph{Separated line classes have disjoint enlargements.}
The fixed-line bound alone does not control the number of possible lines.
We combine it with the separation of their sign classes using the cube's
isoperimetric inequality. For completeness, the
elementary estimate we need is
\begin{equation}\label{eq:cube-boundary}
 |\partial_E S|\ge |S|\log_2(2^r/|S|)
 \qquad(\varnothing\ne S\subseteq\{0,1\}^r).
\end{equation}
Here \(\partial_E S\) consists of cube edges with exactly one endpoint
in \(S\); see also \cite{Harper}.
The bound follows by induction on \(r\), starting with the one-vertex
cube, and with empty contributions interpreted as zero. If the two half-cubes contain
\(a|S|\) and \((1-a)|S|\) members of \(S\), their internal boundary
estimates sum to
\[
 |S|\bigl(r-1-\log_2|S|+h_2(a)\bigr),
\]
where \(h_2\) is binary entropy in bits. At least
\(|1-2a||S|\) edges cross between a member and a nonmember in the two
halves. The inequality
\(h_2(a)\ge1-|1-2a|\), obtained from concavity on each half of
\([0,1]\), proves \eqref{eq:cube-boundary}.
Since an outside vertex has at most \(r\) incident boundary edges,
the outside vertex boundary has size at least
\[
 \frac{|S|}{r}\log_2(2^r/|S|).
\]

For each relevant line \(\ell\), let \(S_\ell\) be its nonempty set
of bad sign vectors. The unique witness line proved above makes these
sets a partition of all bad sign vectors. There are only finitely many such lines, since \(Q\) contains
finitely many lattice points. Distinct \(S_\ell\)'s have Hamming
distance greater than \(dr/4\), by the first part of the proof.
Consequently their Hamming neighborhoods of radius
\(h=\lfloor dr/10\rfloor\) are disjoint.
By \eqref{eq:line-tail}, each \(S_\ell\) has relative size at most
\(e^{-c_0d^2r}\).
As long as an enlarging set has relative size at most
\(e^{-c_0d^2r/2}\), each one-step enlargement increases its size by a
factor at least \(1+c_1d^2\). After \(h\) steps this gives a factor
at least \(e^{c_2d^3r}\).
If the indicated relative-size threshold is reached earlier, the
growth from the initial size is already at least
\(e^{c_0d^2r/2}\), which is an even stronger bound after adjusting
\(c_2\).
The disjoint enlarged sets have total size at most \(2^r\).
Their original total size is therefore at most
\(2^re^{-c_2d^3r}\), proving the lemma.
\end{proof}

\subsection{From differences to endpoint entropy}

The two geometric estimates now give a sampling operation that can be
iterated. We first specify how entropy is averaged over residue maps;
the same convention will be used throughout the sampling schedule.

Fix a batch \(\mathcal P\) of \(k\) distinct primes \(p\equiv1\pmod4\)
in \([T,2T]\), and write
\[
 L_*=\frac1k\sum_{p\in\mathcal P}\log p.
\]
A signed subset \(\mathcal S\) consists of a subset of distinct primes
from \(\mathcal P\), with one of the two conjugate Gaussian factors
chosen for each. Thus its size counts rational primes, and it never
contains both factors of the same prime.
For a finitely supported point variable \(Z\in\Z[i]\), define
\[
 F_{\mathcal S}(Z)=(Z\bmod\pi)_{\pi\in\mathcal S},
 \qquad
 f_Z(s)=\E_{\mathcal S}H(F_{\mathcal S}(Z)),
 \quad 0\le s\le k,
\]
where the \(s\) underlying primes are chosen uniformly without
replacement and their signs independently and fairly.
All index choices are independent of the point variables.
Entropy is computed for each fixed signed subset before averaging.

We have
\[
 f_Z(0)=0,\qquad 0\le f_Z(s)\le sL_*.
\]
The function \(f_Z\) is concave on the integer interval.
Indeed, choose a uniform random ordering of the batch and independent
signs. The difference \(f_Z(s+1)-f_Z(s)\) is the average conditional
entropy of the next coordinate given the first \(s\).
Exchangeability of the unexposed coordinates and the fact that
conditioning decreases entropy show that these differences decrease
with \(s\). This is the usual entropy submodularity argument
\cite{CoverThomas}.

Let \(Z=z_\tau\), where \(\tau\) is any finitely supported random time.
For parameters \(U>0\) and \(0<g<10^{-2}\), put
\(n=\lfloor e^{gU}\rfloor\). For an exact time \(t\), write
\[
 E_t=\{z_t,z_{t+1},\ldots,z_{t+n}\}.
\]
Conditional on \(\tau=t\), choose \(h\) uniformly from the set
\(E_t-E_t\), including zero, and then choose one ordered pair
\((X,Y)\in E_t^2\) with \(X-Y=h\) by a fixed deterministic rule.
Thus differences, not ordered pairs, receive equal weight.
Finally let \(V\) equal \(X\) or \(Y\), each with probability
\(1/2\), using an independent coin. Self-avoidance identifies each
endpoint with a unique walk time in \([t,t+n]\), so this defines a
forward transition of the time law. The transition advances time by
at most \(n\) and uses no prime indices.

\begin{lemma}[Entropy enrichment]\label{lem:enrichment}
Let \(Z=z_\tau\) and \(V\) be the starting point and sampled endpoint
in the preceding transition for a fixed self-avoiding \(D\)-step walk.
Set
\[
 q=\lfloor U/L_*\rfloor,\qquad b=\lfloor g^2U/L_*\rfloor.
\]
Consider any asymptotic family in which \(D\) is fixed,
\(T\to\infty\), \(q\le k\), and
\[
 \frac{g^C U}{\log T}\longrightarrow\infty
 \quad\text{for every fixed }C>0.
\]
Uniformly over the walks, batches, and starting-time laws,
\begin{equation}\label{eq:enrichment}
 \frac{f_V(b)}{bL_*}
 \ge \frac12\left(1+\frac{f_Z(q)}{qL_*}\right)-O_D(g).
\end{equation}
\end{lemma}

The right side retains one half of the starting entropy fraction and
adds one half of the maximal fraction. When successive transitions match
the next input coordinate count to the preceding output count, repetition
therefore drives the entropy deficit downward. Section~\ref{sec:schedule}
will arrange this matching, starting even from a deterministic time.
The error \(O_D(g)\) is uniform; choosing the scale large enough permits
\(g\) to decrease as well.

\begin{proof}
The regime ensures \(1\le b<q\) for sufficiently large parameters.
We first show that a fresh difference coordinate has nearly maximal
entropy even after the endpoint prefixes are known. We then use the
entropy chain rule to retain the contribution from \(Z\).

\paragraph{Entropy of one further difference coordinate.}
Choose a uniform signed set \(\mathcal Q\) of size \(q\), and within it
a uniform signed prefix \(\mathcal S\) of size \(b\). Write
\[
 B_x=F_{\mathcal S}(X),\qquad B_y=F_{\mathcal S}(Y),\qquad h=X-Y.
\]
If \(j\) is a uniform remaining signed coordinate and
\(h_j=F_j(h)\), we first prove
\begin{equation}\label{eq:difference-coordinate}
 \E_{\mathcal Q,\mathcal S,j}H(h_j\mid B_x,B_y)
 \ge (1-O_D(g))L_*.
\end{equation}

Fix an exact base time \(\tau=t\) with positive probability, and use \(A\) from
Lemma~\ref{lem:differences} for its segment. That lemma gives
\[
 gU-O_D(1)\le\log A\le2gU+O_D(1).
\]
Put
\[
 r=\left\lceil(1+10g)\frac{\log A}{L_*}\right\rceil.
\]
For sufficiently large parameters,
\(gq/2\le r\le3gq\le q-b\).
Choose a uniform size-\(r\) subset \(\mathcal J\) of the remaining
coordinates. Averaging over \(\mathcal Q,\mathcal S\), its underlying
primes form a uniform sample of size \(r\) from the batch.
Their log-product has mean \(rL_*\) and variance \(O(r)\), since
the values \(\log p\) lie in an interval of length \(\log2\) and
sampling is without replacement. Hence
\begin{equation}\label{eq:product-concentration}
 \Pr\left\{\sum_{p\in\mathcal J}\log p
                 <(1+5g)\log A\right\}
 =O\left(\frac1{g^2rL_*^2}\right).
\end{equation}

In coordinates along the diameter rectangle, differences of
differences from the segment lie in an origin-centered rectangle
of half-lengths \(4R,4W\), whose product is \(16A\).
On the complementary event in \eqref{eq:product-concentration},
writing \(P=\prod_{p\in\mathcal J}p\), we have
\[
 P^{1-g}\ge A^{(1+5g)(1-g)}\ge A^{1+3g}>16A
\]
for sufficiently large parameters. Lemma~\ref{lem:rectangle},
with \(d=g\), shows that reduction in the coordinates
\(\mathcal J\) is injective on the difference set, except with
additional probability \(O(e^{-cg^3r})\).
Let
\[
 \epsilon=O\left(\frac1{g^2rL_*^2}+e^{-cg^3r}\right)
\]
be the total exceptional probability.

Here all errors are uniform in the exact base time.
Indeed \(r\) is bounded above and below by constant multiples of
\(gU/L_*\), while \(L_*\) is comparable to \(\log T\).
The assumed regime implies
\[
 \frac1{gq}=o(g),\qquad
 \frac1{g^2rL_*^2}=o(g),\qquad
 g^3r/\log(1/g)\longrightarrow\infty
\]
when \(g\to0\); for \(g\) bounded away from zero the last assertion
is replaced simply by \(g^3r\to\infty\).
For example, the condition with \(C=5\) makes
\(g^3r\) larger than an unbounded multiple of \(1/g\).
Thus \(\epsilon=o(g)\), ceiling errors are \(o(g)\) relatively,
and \(g\log A\to\infty\), as used above.

On a good choice of \(\mathcal J\), uniformity on the distinct
differences and injectivity give
\[
 H(F_{\mathcal J}(h)\mid\tau=t)\ge\log A-O_D(1).
\]
Conditioning additionally on \(B_x,B_y\) costs at most
\(2b\log(2T)\). Averaging the fixed-index entropies therefore yields
\[
 \E_{\mathcal Q,\mathcal S,\mathcal J}
 H(F_{\mathcal J}(h)\mid B_x,B_y,\tau=t)
 \ge (1-\epsilon)(\log A-O_D(1))-2b\log(2T).
\]
For every fixed index choice, conditional subadditivity gives
\[
 H(F_{\mathcal J}(h)\mid B_x,B_y,\tau=t)
 \le\sum_{j\in\mathcal J}H(h_j\mid B_x,B_y,\tau=t).
\]
Consequently,
\[
 \frac1r\E_{\mathcal Q,\mathcal S,\mathcal J}
 H(F_{\mathcal J}(h)\mid B_x,B_y,\tau=t)
 \le \E_{\mathcal Q,\mathcal S,j}
 H(h_j\mid B_x,B_y,\tau=t),
\]
where \(j\) is uniform among the remaining coordinates.
This remains true although \(r\) depends on \(t\): at this stage
\(t\) is fixed, and every member of the uniform \(\mathcal J\)
has the same uniform remaining-coordinate marginal.
Since
\[
 \frac{\log A}{rL_*}\ge1-O(g),\qquad \frac br=O(g),
\]
we obtain the right side of \eqref{eq:difference-coordinate},
still with the extra conditioning \(\tau=t\).
Averaging over \(t\), and then removing the conditioning on \(\tau\),
proves \eqref{eq:difference-coordinate}.

\paragraph{Retaining the starting entropy.}
We have established near-maximal entropy for a further coordinate of
\(h=X-Y\), after revealing the endpoint prefixes.
It remains to turn that information into endpoint entropy while keeping
the contribution from the starting law.
For fixed index choices, the pair
\((F_j(X),F_j(Y))\) is equivalent to \((h_j,F_j(X))\).
The entropy chain rule and conditioning inequalities give
\begin{align}
 H(h_j\mid B_x,B_y)
 &=H(F_j(X),F_j(Y)\mid B_x,B_y)
       -H(F_j(X)\mid B_x,B_y,h_j) \notag\\
 &\le H(F_j(X)\mid B_x)+H(F_j(Y)\mid B_y)
       -H(F_j(X)\mid B_x,B_y,h).
 \label{eq:enrichment-chain}
\end{align}
In particular the last conditioning is on the entire difference
\(h\), not just on \(h_j\); using it decreases the subtracted
entropy and preserves the displayed inequality.
Summing these subtracted entropies over the \(q-b\) remaining
coordinates gives at least
\[
 H(F_{\mathcal Q}(X)\mid B_x,B_y,h),
\]
because the prefix coordinates are already known from \(B_x\).
Its average is at least
\[
 f_X(q)-2b\log(2T)-H(h).
\]
Also \(F_{\mathcal Q}(Z)\) is determined by
\(F_{\mathcal Q}(X)\) and \(X-Z\), so
\[
 f_X(q)\ge f_Z(q)-H(X-Z).
\]
Both \(h\) and \(X-Z\) are lattice displacements of length at most
\(Dn\). The counting bound of Lemma~\ref{lem:displacement} gives
\[
 H(h)+H(X-Z)\le4\log(n+1)+O_D(1)=O_D(gU).
\]
Consequently the average contribution of the subtracted term in
\eqref{eq:enrichment-chain} is at least
\[
 \frac{f_Z(q)-O_D(gU)}{q-b}
 \ge \frac{f_Z(q)}q-O_D(g)L_*.
\]

The averages of the first two terms on the right side of
\eqref{eq:enrichment-chain} are, respectively,
\(f_X(b+1)-f_X(b)\) and \(f_Y(b+1)-f_Y(b)\).
Combining with \eqref{eq:difference-coordinate} proves
\[
 [f_X(b+1)-f_X(b)]+[f_Y(b+1)-f_Y(b)]
 \ge L_*+\frac{f_Z(q)}q-O_D(g)L_*.
\]
Concavity and \(f_X(0)=f_Y(0)=0\) bound these forward differences
above by \(f_X(b)/b\) and \(f_Y(b)/b\).
Finally, concavity of entropy under a mixture gives
\(f_V(b)\ge\frac12(f_X(b)+f_Y(b))\).
Dividing by \(bL_*\) proves \eqref{eq:enrichment}.
\end{proof}
